\section{Introduction}\label{dp:sec:intro}

Let $K$ be a nonarchimedean nontrivially valued field, fix $\varpi\in K$ with $0<\abs\varpi<1$, and let $P,W$ be normed $K$-spaces, not necessarily complete and not necessarily ultrametric. A \emph{bounded polynomial law} of degree $n$, or \emph{law}, is a map $a\colon P\to W$ that is a homogeneous polynomial of degree $n$ on every finite-dimensional subspace and satisfies $\norm{a(h)}\le C\norm h^n$ for some real $C$; the least such $C$ is $\norm a$. A map $f$ from an open subset of $P$ to $W$ is \emph{divided-power analytic}, or \emph{DP}, at $x_0$ if $f(x_0+h)=\sum_{n\ge0}a_n(h)$ for $\norm h<R$, with laws $a_n$ of degree $n$ and $\norm{a_n}\le MR^{-n}$ (Definition~\ref{dp:def:laws}). The name does not mean a series $\sum_na_nx^n/n!$. No division occurs, and if $K$ is complete and $P$ is finite-dimensional, a DP germ is a convergent power series (Theorem~\ref{dp:thm:conservativity}(5)). The name refers to the target of the universal law: a law of degree $n$ is a bounded linear map on the span $\Gamma^n(P)$ of the powers $h^{\otimes n}$, for the largest seminorm $\norm\cdot_\Gamma$ with $\norm{h^{\otimes n}}_\Gamma\le\norm h^n$ (Theorem~\ref{dp:thm:normed-roby}). Without the bound, this is Roby's description of laws by divided powers~\cite[Th.~IV.1, Prop.~IV.5]{Ro}.

A map is \emph{analytic} at $x_0$ if each $a_n$ above is the diagonal of an $n$-linear map $\beta_n\colon P^n\to W$ with $\norm{\beta_n}\le MR^{-n}$. Analytic germs are DP, with the same radius and bound. Roby rejected multilinear diagonals as a definition of polynomial maps: factorials get in the way, even for vector spaces over a finite field, and the multilinear map is not unique~\cite[p.~214]{Ro}. Section~\ref{dp:sec:laws} identifies laws with Roby's homogeneous polynomial laws. So, without the bounds, the laws of positive degree of a DP germ form a formal law in Roby's sense~\cite[Th.~VI.1, Th.~VI.2]{Ro}. So a DP germ is a Roby formal law plus a constant, with geometric bounds. The algebra of composition and of the formal inverse is Roby's~\cite[VI~\S8, Prop.~VI.6, Cor.; Th.~VI.3, Prop.~VI.8]{Ro}. We add only the estimates that make the series converge (Theorems~\ref{dp:thm:calculus} and~\ref{dp:thm:inverse}).

\subsection*{Laws that are not ambient}

Paper~\cite{I} enlarges the target. A map $f$ is \emph{ambient analytic}, or \emph{AA}, at $x_0$ if $j\circ f$ is analytic at $x_0$ for a linear isometry $j$ of $W$ into a normed space, with the diagonals of the coefficients in $j(W)$. A law $a$ of degree $n$ is \emph{ambient} if $j\circ a$ is the diagonal of a bounded $n$-linear map for such a $j$. An AA germ is a DP germ whose laws are ambient, with geometric bounds on the norms of these multilinear maps~\cite[Proposition~14.9]{I}. Over every $K$ of characteristic $p>0$, the class AA is not closed under composition~\cite[Theorem~14.17]{I}. If $\operatorname{char}K=p>0$, then in degrees below $p$ every law is the diagonal of a bounded multilinear map (Section~\ref{dp:sec:nonambient}). From degree $p$ on this can fail, even after an enlargement of the target.

\begin{introtheorem}[laws that are not ambient]\label{dp:intro:A}
Let $\operatorname{char}K=p>0$.
\begin{enumerate}
\item Over every $K$, in every degree $n\ge p$, there is a law of degree $n$ and norm at most $1$ that is not ambient: the map $(\lambda,h)\mapsto\lambda^{n-p}h^{\otimes p}$, with $\lambda^0=1$, from $K\oplus\mathcal E$, with the maximum norm, to $\Gamma^p(\mathcal E)$ with the norm $\norm\cdot_\Gamma$, for a suitable $\ell^1$-sum $\mathcal E$ of finite-dimensional spaces.
\item If $K$ is discretely valued or $p=2$, composition fails on $H_0=c_0\oplus G$, with the maximum norm, where $c_0=c_0(\N,K)$ and $G$ is a suitable $\ell^1$-sum of finite-dimensional spaces. Let $X=c_0$, and give $\Gamma^p(X)$ the projective norm. The map $f_X(x,w)=(x^{\otimes p},w)$ into $\Gamma^p(X)\oplus G$ is AA at every point, the map $g_X(s,w)=s\otimes w$ into the projective tensor product of $\Gamma^p(X)$ and $G$ is the diagonal of a bounded bilinear map, and $g_X\circ f_X$ is a law of degree $p+1$, DP everywhere and AA at no point.
\end{enumerate}
\end{introtheorem}

This is Proposition~\ref{dp:prop:degree-p}(2) and Theorem~\ref{dp:thm:c0-failure}. Over densely valued $K$ of odd characteristic, whether this failure occurs on $c_0\oplus G$ is open~\cite[Remark~14.20]{I}.

\subsection*{The two descents and minimality}

Consider classes of germs with Banach values and three properties of such a class. (C1): it contains the analytic germs. (C2), \emph{target descent}: if $j\colon W\to Z$ is a linear isometry of Banach spaces and $j\circ f$ is in the class, then $f$ is. (C3), \emph{source descent}: if $Q\colon P'\to P$ is a bounded linear surjection with a preimage constant and $f\circ Q$ is in the class, then $f$ is. Here $Q$ has \emph{preimage constant} $\mathsf c$ if every $h\in P$ has a preimage of norm at most $\mathsf c\norm h$.

\begin{introtheorem}[the smallest class with these descents]\label{dp:intro:B}
\leavevmode
\begin{enumerate}
\item The DP germs with Banach values are the smallest class with \textup{(C1), (C2), (C3)}.
\item Among germs with ultrametric sources and ultrametric targets, the DP germs are the smallest class with \textup{(C1)} and \textup{(C3)}.
\end{enumerate}
\end{introtheorem}

This is Theorem~\ref{dp:thm:minimal}(3) and Theorem~\ref{dp:thm:ultrametric-minimal}(2). Item~(2) is where ultrametric norms simplify. There the targets need not be complete. No target descent is needed: $P$ is the image of $c_{00}(\mathsf I)$, the finitely supported functions on a set $\mathsf I$ with the supremum norm, under a contractive surjection with a preimage constant, and a law on $c_{00}(\mathsf I)$ with values in an ultrametric space is already the diagonal of a multilinear map, with constant $\abs\varpi^{-n}$. Analytic germs have (C1) and, in every characteristic, fail (C2). In characteristic $p$ they fail (C3), and AA, the smallest class with (C1) and (C2), fails it too (Proposition~\ref{dp:prop:other-classes}).

Minimality is a statement about these axioms, and on general sources it depends on target descent into non-ultrametric spaces. Fix $\mathsf e_n\in K$ with $n^n\le\abs{\mathsf e_n}<\abs\varpi^{-1}n^n$, and put
\[
 \mathsf f(h)=\sum_{n\ge1}\mathsf e_nh_1h_2\cdots h_n
\]
on the open unit ball of $\ell^1(\N,K)$, with values in $\Kh$. By the arithmetic--geometric mean inequality, the laws have norm below $\abs\varpi^{-1}$. So $\mathsf f$ is DP at $0$. It is AA, so every class with (C1), (C2) contains it. In every characteristic it is not analytic, since a multilinear map whose diagonal is the law of degree $n$ has norm at least $n^n$ (Example~\ref{dp:ex:l1}). The series converges because of an inequality between real numbers on the source, and no divided power is involved.

Let CA be the class of germs \emph{analytic after a cover}: germs $f$ with $f\circ Q$ analytic at a preimage of the base point, for some bounded linear surjection $Q$ with a preimage constant. Among germs with Banach values it is the smallest class with (C1) and (C3) (Lemma~\ref{dp:lem:cover-flat}(3)). For $K$ complete it also has composition, pairing, inverse and implicit functions, it contains every law, and it does not contain $\mathsf f$ (Theorem~\ref{dp:thm:descents-independent}(1)). So on general sources (C1) and (C3) reach only CA. A class with (C1) and (C3) whose germs may have non-ultrametric sources contains every DP germ on an ultrametric source, with every target (Corollary~\ref{dp:cor:um-cover}); the cover used there is not ultrametric. In characteristic $p$, whether, among germs with ultrametric sources and Banach targets, the DP germs are the smallest class with (C1), (C2), (C3) is open (Question~\ref{dp:q:c0}(2)). The other way up to DP on general sources is an axiom for classes closed under translation in the source, the axiom (S): if laws $b_n$ of degree $n$ have germs at $0$ in the class and geometric bounds on their norms, then the germ of $\sum_nb_n$ at $0$ is in the class. Such a class with (C1), (C3), (S) contains every DP germ with Banach values, and in characteristic $0$ (C1) and (S) suffice (Propositions~\ref{dp:prop:s-axiom} and~\ref{dp:prop:s-char0}). The axiom (S) is close to the definition of DP.

\subsection*{Composition is not enough}

The failure of composition in Theorem~\ref{dp:intro:A} does not lead to DP. In characteristic $p$, repairing composition closes up strictly below DP. Let $\mathrm{AA}^\circ$ be the smallest class of germs with Banach values that contains the analytic germs and is closed under composition, pairing and target descent. Let $\pi$ be the projective seminorm on $\Gamma^N(P)\subseteq P^{\otimes N}$. A germ is AA if and only if the linear maps on $\Gamma^N(P)$ given by its laws have geometric bounds for $\pi$, and DP if and only if they have them for $\norm\cdot_\Gamma$. Between the two lie the \emph{tree norms} $\mathfrak t_N$, iterated projective seminorms along the compositions of the integer $N$, with $\pi\le\mathfrak t_N\le\norm\cdot_\Gamma$. A germ is \emph{tree analytic}, or TA, if its linear maps have geometric bounds for $\mathfrak t_N$ (Definitions~\ref{dp:def:tree} and~\ref{dp:def:ta}). Let $\cC_p$ be the class of DP germs whose laws of degree at most $p$ are ambient. Section~\ref{dp:sec:chain} also has two variants of $\cC_p$: a local one $\cC_p^{\rm loc}$, the germs in $\cC_p$ at every point near the base point, and a uniform one $\cC_p^{\rm u}$.

\begin{introtheorem}[the chain]\label{dp:intro:C}
Let $\operatorname{char}K=p>0$, with $K$ complete or not.
\begin{enumerate}
\item Among germs with Banach values,
\[
 \mathrm{AA}\subsetneq\mathrm{AA}^\circ\subseteq\mathrm{TA}\subseteq\cC_p^{\rm u}\subseteq\cC_p^{\rm loc}\subsetneq\cC_p\subsetneq\mathrm{DP}.
\]
\item TA and $\cC_p$ are closed under composition and pairing, and between Banach spaces they have inverse functions.
\item The map $\gamma^\sharp_p(\lambda,h)=h^{\otimes p}$ on $K\oplus\mathcal E$, for the space $\mathcal E$ of Theorem~\ref{dp:intro:A} and with the norm $\norm\cdot_\Gamma$ on values, is DP at every point, and its germ at no point is in $\cC_p$.
\item TA and $\cC_p$ fail \textup{(C3)}. A class with \textup{(C1), (C2)} has \textup{(C3)} only if it contains every DP germ with Banach values.
\end{enumerate}
\end{introtheorem}

This is Theorem~\ref{dp:thm:chain}, with Theorems~\ref{dp:thm:ta-composition}, \ref{dp:thm:ta-inverse}, \ref{dp:thm:cp}, \ref{dp:thm:cp-inverse} and~\ref{dp:thm:singles-out}. By Theorem~\ref{dp:thm:singles-out}, TA has (C1), (C2), composition, pairing, inverse and implicit functions, re-expansion on the whole ball, derivatives and bundle lifting. Of the properties listed there, only (C3) is not shared by TA, and with (C1), (C2) it forces DP. The separations in Theorem~\ref{dp:intro:C} and the failure of (C3) are all in characteristic $p$. In characteristic $0$, whether $\mathrm{AA}\ne\mathrm{DP}$ and whether $\mathrm{TA}\ne\mathrm{DP}$ are open. Between ultrametric spaces the two papers study one class, from outside in~\cite{I} and from inside here: there $\mathrm{AA}=\mathrm{TA}=\mathrm{DP}$.

\subsection*{Jets as a second test}

Jets are not on that list, and they test a class in another way. For DP manifolds $\mathcal M,\mathcal N$ with normed models, in every characteristic, the jet space $J^r(\mathcal M,\mathcal N)$ is a DP manifold, with DP forgetful maps, prolongations and functoriality (Theorem~\ref{dp:thm:dp-jets}). For analytic manifolds the question is governed by a section. We say that $(\mathrm{Sec}_k)$ holds for $(P,W)$ if the diagonal map from bounded $k$-linear maps $P^k\to W$ to laws $P\to W$ of degree $k$ has a bounded linear section. The \emph{operator family} of a chart change sends a point to the linear map on jets with values in $W$ that the chart change induces at that point. If $(\mathrm{Sec}_k)$ holds for $1\le k\le r$, the jet bundles of analytic manifolds with models $P,W$ have analytic transitions and operator families. The \emph{shear manifold} of a normed space $V$ is $V\times V\times K$ with the two global charts $\id$ and $(x,z,\lambda)\mapsto(x,z+\lambda x,\lambda)$. If on the shear manifold of $V$ the operator family of the jet bundle with values in $W$ is analytic at one point, then $(\mathrm{Sec}_k)$ holds for $(V,W)$ and $1\le k\le r$ (Theorem~\ref{dp:thm:analytic-jets}).

Let $K$ be complete of characteristic $p>0$ and $r\ge p$. The scalar jet bundles $J^r(\mathcal M;K)$ of all analytic manifolds $\mathcal M$ with ultrametric models have analytic transitions and operator families if and only if $K$ is spherically complete (Corollary~\ref{dp:cor:jet-dichotomy}). This separates analyticity from $\mathrm{AA}=\mathrm{DP}$. It does not separate AA from DP, as the family that fails is a map between ultrametric spaces, where AA and DP coincide. The separation from AA is Proposition~\ref{dp:prop:jets-not-ambient}(4), over every $K$ of characteristic $p$: for $r\ge p$, on the shear manifold of $\mathcal E$, with target the completion of $\Gamma^p(\mathcal E)$ for $\norm\cdot_\Gamma$, the operator family of the jet bundle with values in that target is DP at every point and misses $\cC_p$ at every point. So it is neither AA nor TA at any point.

\subsection*{The whole picture}

Between ultrametric spaces the classes above, apart from the analytic germs, are one class. The AA germs, the TA germs, the strongly DP germs of Section~\ref{dp:sec:strong}, the CA germs and the DP germs coincide there (\cite[Theorem~14.39(1)]{I}, Theorem~\ref{dp:thm:singles-out}(4), Lemmas~\ref{dp:lem:strong-formal}(3) and~\ref{dp:lem:cover-flat}(5)). The comparison with analyticity then turns on the field. For $k\ge1$, normed spaces $E,D,F,F'$, $u\in\cL(E,D)$, $v\in\cL(F,F')$ and $k$-linear maps $\mathsf m\colon D^k\to F$, let $B(u,v)(\mathsf m)=v\circ\mathsf m\circ(u,\dots,u)$ be the joint action, and let $A^k_{E,D;F}(u)=B(u,\id_F)$ on alternating $\mathsf m$. For Banach spaces this is the formula of~\cite[n\textsuperscript{o}~7.8.1]{Bou} for $\alpha_k(\id_F,u)$~\cite[Section~10.2]{I}. Bourbaki defines the functor $\alpha_k$ when $K$ has characteristic $0$ or the bundles have finite rank (Section~\ref{dp:sec:functors}); we use the formula for all normed spaces. Let $K$ be complete, and consider two statements.
\begin{itemize}
\item[$(\mathrm D_K)$] Every map between open subsets of ultrametric Banach spaces that is DP at a point is analytic there.
\item[$(\mathrm B_K)$] For all ultrametric Banach spaces $E,D,F,F'$ and every $k\ge1$, the maps $A^k_{E,D;F}$ and $B$, restricted to alternating maps, are analytic everywhere.
\end{itemize}

\begin{introtheorem}[agreement with analyticity]\label{dp:intro:D}
Let $K$ be complete.
\begin{enumerate}
\item $(\mathrm D_K)$ implies $(\mathrm B_K)$. Both hold if $\operatorname{char}K=0$ or if $K$ is discretely valued. Both fail if $\operatorname{char}K=p>0$ and $K$ is not spherically complete.
\item Let $P$ be ultrametric and $W$ an ultrametric Banach space. If $f$ is DP at $x_0$ and not analytic there, then $\operatorname{char}K=p>0$, $K$ is densely valued, the completion of $P$ is not linearly homeomorphic to any $c_0(\mathsf I,K)$, and $W$ has no equivalent ultrametric spherically complete norm.
\item Let $P$ be finite-dimensional and $W$ any normed space. A map from an open subset of $P$ to $W$ is DP at a point if and only if it is analytic there.
\end{enumerate}
\end{introtheorem}

This is Theorem~\ref{dp:thm:field-level} and Theorem~\ref{dp:thm:conservativity}(2),(5). So, wherever $(\mathrm B_K)$ is decided here, $(\mathrm D_K)$ is decided with the same answer. Between ultrametric Banach spaces, the disagreements listed here occur over the complete, not spherically complete $K$ of characteristic $p$, for some spaces, and there, for $k\ge p$ and some ultrametric Banach spaces $E,D$, the map $A^k_{E,D;K}$ is analytic at no point (Proposition~\ref{dp:prop:strict}(1)). One case is undecided. If $K$ is densely valued and spherically complete of characteristic $p$, we decide neither $(\mathrm D_K)$ nor $(\mathrm B_K)$. Whether $(\mathrm B_K)$ implies $(\mathrm D_K)$ is open (Section~\ref{dp:sec:questions}).

Section~\ref{dp:sec:ledger} lists the disagreements between DP and analyticity stated in this paper and in~\cite{I}, by field and by norms (Table~\ref{dp:tab:ledger}), and says what forces each one (Theorem~\ref{dp:thm:forcing}). We do not claim that the list is complete: in characteristic $p$, on ultrametric sources with non-ultrametric targets, whether DP and AA differ is open.

The strong norm appears both as a norm and as a competing class. For a law $a$ and vectors $y_1,\dots,y_l\in P$, expand $a(\sum_it_iy_i)$ as a polynomial in scalars $t_1,\dots,t_l$. The \emph{strong norm} of $a$ is the least constant $C$ such that the coefficient of each monomial $\prod_it_i^{d_i}$ has norm at most $C\prod_i\norm{y_i}^{d_i}$, for all $l$ and all $y_1,\dots,y_l$, and a germ is \emph{strongly DP} if its laws have geometric bounds in this norm. As a norm it loses no radius on ultrametric spaces (Theorems~\ref{dp:thm:strong-inverse} and~\ref{dp:thm:strong-recenter}). For $W$ an ultrametric Banach space and every $P$, a germ is strongly DP if and only if it is analytic after a cover, if and only if its laws descend to the ultrametric hull of $P$, the quotient of $P$ by the kernel of the seminorm $\mathsf u$ defined below, with geometric bounds there (Theorem~\ref{dp:thm:strong-um}). As a class it keeps (C2) and (C3), and it can lose (C1) (Theorem~\ref{dp:thm:strong-descent}). On $\ell^1(\N,K)$, with a non-ultrametric target, an analytic germ need not be strongly DP (Example~\ref{dp:ex:an-not-sdp}). On the same source, over every $K$ and in every characteristic, strongly DP germs have neither composition nor inverse functions (Theorem~\ref{dp:thm:strong-failure}).

We also have two results for all norms. If $W$ is complete, a map that is DP at $x_0$ with radius $R$ is DP at every point of the open ball of radius $R$ around $x_0$, with no hypothesis on $K$ or on the norms (Theorem~\ref{dp:thm:whole-ball}), though over some discretely valued $K$ its largest radius at a point $x_1$ of that ball can be smaller than $R-\norm{x_1-x_0}$ (Example~\ref{dp:ex:radius-drop}). Every law $a$ of degree $n$ extends to a law $a_{K'}$ on the base change to every valued extension $K'$ of $K$, with $\norm{a_{K'}}\le C_0^n\norm a$ for $C_0=8\mathrm e/\abs\varpi^2$ (Theorem~\ref{dp:thm:base-change-law}). Base change does not separate DP from analyticity by stability: both classes are stable, analytic germs with no loss of radius (Theorem~\ref{dp:thm:an-base-change}(1)), and DP germs with a possible loss of radius (Theorem~\ref{dp:thm:germ-base-change}(1)), which occurs with the maximum norm on the base change (Example~\ref{dp:ex:base-radius}). A DP germ whose base change is analytic or AA over $K'$ is already AA over $K$, since base change is a particular enlargement of the target (Theorems~\ref{dp:thm:an-base-change}(3) and~\ref{dp:thm:aa-base-change}(3)).

\subsection*{Organization and norms}

Section~\ref{dp:sec:conv} fixes conventions and recalls what we use from~\cite{I}. Sections~\ref{dp:sec:laws} and~\ref{dp:sec:germs} give the calculus of laws and germs. Section~\ref{dp:sec:nonambient} gives laws that are not ambient, with proofs in Appendix~\ref{dp:app:proofs}. Section~\ref{dp:sec:descent} gives descent and minimality, and Section~\ref{dp:sec:chain} the classes between AA and DP. Sections~\ref{dp:sec:forms}, \ref{dp:sec:quotients} and~\ref{dp:sec:base-change} treat forms, difference quotients and base change. Section~\ref{dp:sec:agreement} compares DP with analyticity, Section~\ref{dp:sec:functors} treats vector functors and n\textsuperscript{o}~7.8 of~\cite{Bou}, and Section~\ref{dp:sec:jets} treats jets. Sections~\ref{dp:sec:strong} and~\ref{dp:sec:cover} study the strong norm and CA, Section~\ref{dp:sec:ledger} is the list of disagreements, and Section~\ref{dp:sec:questions} collects questions.

We work with general norms as far as they extend naturally. The mixed-term bound, the whole-ball theorem, composition, inverse functions, descents, bundles, forms, jets and base change need no ultrametric hypothesis. Ultrametric norms enter where they are used, and we say so there. The main places and reasons are these.
\begin{itemize}
\item Theorem~\ref{dp:thm:mixed}(4). On an ultrametric source the mixed-term bound loses no factor $n^n$. With an ultrametric target as well, Lemma~\ref{dp:lem:strong-coefficients} gives the constants.
\item Theorem~\ref{dp:thm:ultrametric-minimal}. Between ultrametric spaces, source descent reaches DP without target descent, since laws on the cover are diagonals with geometric constants.
\item Sections~\ref{dp:sec:agreement} and~\ref{dp:sec:strong}, on targets. In an ultrametric target a sum is bounded by its maximum. So multilinear maps into it see only the hull seminorm $\mathsf u(x)=\inf\max_i\norm{x_i}$ of the source, the infimum over the decompositions $x=\sum_ix_i$, and a coefficient, a sum with one term for each ordering, is no larger than its largest term.
\end{itemize}
Every theorem is proved here or cited from~\cite{I} or from the bibliography.
